\documentclass[11pt]{article}
\usepackage[margin=1in]{geometry}
\usepackage{amsmath,amssymb}
\usepackage[T1]{fontenc}
\usepackage[utf8]{inputenc}
\usepackage{parskip}

\begin{document}

\title{Technical Review of ``Using Double-Entry Accounting to Program Quantum Computers''}
\author{}
\date{}
\maketitle

This report refers to paragraph numbers from an extracted version of the DOCX because the Word file does not have stable line numbers.

\section*{Recommendation}

The paper should be treated as a \textbf{major revision}. Its constructive core is real: the algebra for the basic gates is mostly correct, and the stated Toffoli decomposition checks out. However, several surrounding claims are stronger than the mathematics supports, and two foundational statements are incorrect as written.

\section*{Summary}

The strongest part of the paper is the explicit symbolic encoding of phase and bit-flip operations and the step-by-step accounting version of a Toffoli circuit. I independently checked the decomposition
\[
\mathrm{Toffoli}
=
H_3 X_{23} T_3^\dagger X_{13} T_3 X_{23} T_3^\dagger X_{13} T_3
T_2^\dagger X_{12} H_3 T_2^\dagger X_{12} T_1 S_2,
\]
and it matches the standard Toffoli action on all eight computational-basis states.

The main weaknesses are conceptual and foundational. In particular, the paper currently misstates the composition postulate, incorrectly claims that \(X\), \(Z\), and \(T\) generate all single-qubit unitaries, gives an incomplete derivation of \(H^2=I\), and overstates the role of reversibility in distinguishing quantum from classical computation.

\section*{Major Revisions}

\subsection*{1. Postulate 3 is incorrect as written}

The paper says that ``the state of a composite system is the tensor product of the states of the component systems.'' That is false in quantum mechanics. The correct statement is that the \emph{state space} of the composite system is the tensor product of the component state spaces.

For two qubits, the correct Hilbert-space statement is
\[
\mathcal H_{12}=\mathcal H_1\otimes \mathcal H_2 \cong \mathbb C^2 \otimes \mathbb C^2.
\]
A general two-qubit pure state has the form
\[
|\psi\rangle=\alpha_{00}|00\rangle+\alpha_{01}|01\rangle+\alpha_{10}|10\rangle+\alpha_{11}|11\rangle,
\]
and in general it does \emph{not} factor as
\[
|\psi_1\rangle \otimes |\psi_2\rangle.
\]

A standard counterexample is
\[
|\Phi^+\rangle=\frac{|00\rangle+|11\rangle}{\sqrt2}.
\]
If one tries to write
\[
|\Phi^+\rangle=(a|0\rangle+b|1\rangle)\otimes(c|0\rangle+d|1\rangle),
\]
then comparing coefficients gives
\[
ac=\frac{1}{\sqrt2}, \qquad bd=\frac{1}{\sqrt2}, \qquad ad=0, \qquad bc=0,
\]
which is impossible. As written, the paper's statement excludes entangled states, which are essential to the theory.

\paragraph{Suggested replacement text.}
``The state space of a composite system is the tensor product of the state spaces of its component systems. Computational basis states are product states, but general states may be superpositions of product states and need not factor.''

\subsection*{2. The claim that all single-qubit unitaries can be built from \(X\), \(Z\), and \(T\) is false}

The paper states that once \(X\), \(Z\), and \(T\) are available, all other single-qubit unitary operations can be constructed. That is not correct.

The gates
\[
X=\begin{bmatrix}0&1\\1&0\end{bmatrix}, \qquad
Z=\begin{bmatrix}1&0\\0&-1\end{bmatrix}, \qquad
T=\begin{bmatrix}1&0\\0&e^{i\pi/4}\end{bmatrix}
\]
generate only basis permutations together with phases. Any product of these gates is a monomial unitary, i.e. it has exactly one nonzero entry in each row and each column. Equivalently, every such product has one of the two forms
\[
\begin{bmatrix}
e^{i\alpha} & 0\\
0 & e^{i\beta}
\end{bmatrix}
\qquad \text{or} \qquad
\begin{bmatrix}
0 & e^{i\alpha}\\
e^{i\beta} & 0
\end{bmatrix}.
\]
Such gates never map a basis state to a nontrivial superposition.

But
\[
H|0\rangle=\frac{|0\rangle+|1\rangle}{\sqrt2},
\]
so \(H\) cannot be generated by \(X\), \(Z\), and \(T\) alone.

\paragraph{Suggested replacement text.}
``The gates \(X\), \(Z\), and \(T\) generate a useful family of bit-flip and phase-shift operations. To generate superpositions, one must add a gate such as \(H\).''

\subsection*{3. The proof that \(H^{-1}=H\) is incomplete}

The paper correctly writes
\[
H=\frac{X+Z}{\sqrt2}.
\]
That identity is fine. The problem is the explanation that follows. The paper appears to justify \(H^{-1}=H\) only from the facts that \(X\) and \(Z\) are self-inverse. That is not enough.

The right calculation is
\[
H^2=\frac12(X+Z)^2
=\frac12\left(X^2+XZ+ZX+Z^2\right).
\]
Now one uses
\[
X^2=I, \qquad Z^2=I, \qquad XZ=-ZX.
\]
The anticommutation relation is essential because it gives
\[
XZ+ZX=0.
\]
Hence
\[
H^2=\frac12(2I)=I.
\]
The conclusion is correct, but the derivation is incomplete unless the paper explicitly invokes \(XZ=-ZX\).

\paragraph{Suggested replacement text.}
``Because \(H=(X+Z)/\sqrt2\), we have \(H^2=\tfrac12(X^2+XZ+ZX+Z^2)\). Using \(X^2=Z^2=I\) and \(XZ=-ZX\), it follows that \(H^2=I\), so \(H^{-1}=H\).''

\subsection*{4. The CNOT example is internally inconsistent}

The paper says the controlled-\(X\) gate changes \(|10\rangle\) into \(|11\rangle\), which is correct for control qubit 1 and target qubit 2:
\[
X_{12}|10\rangle=|11\rangle.
\]
But the accounting example shown immediately after that corresponds to
\[
|11\rangle \to |10\rangle,
\]
not \(|10\rangle \to |11\rangle\).

Given the paper's earlier convention that \(EA\) denotes \(|01\rangle\), the string \(AA\) denotes \(|11\rangle\) and \(AE\) denotes \(|10\rangle\). So the notation
\[
AA \xrightarrow{X_{12}} AE
\]
is a correct example of CNOT, but it is \emph{not} the same example as \(|10\rangle \to |11\rangle\).

\paragraph{Suggested correction.}
Either write
\[
AE \xrightarrow{X_{12}} AA
\]
to match \(|10\rangle \to |11\rangle\), or change the prose so the example is explicitly \(|11\rangle \to |10\rangle\).

\subsection*{5. The paper overstates what reversibility shows}

Several passages suggest that reversibility is what distinguishes quantum from classical computation, or that Toffoli ``distinguishes'' quantum from classical computation. That is not right.

Toffoli is a \emph{classical reversible} gate. Reversible classical computation has been studied for decades. The fact that a gate is reversible does not make it quantum. What is specifically quantum is the combination of
\[
\text{superposition} + \text{interference} + \text{entanglement} + \text{unitary evolution}.
\]

Landauer's principle is also being used too broadly. Landauer's principle concerns logically irreversible operations, not classical computation as such. A classical reversible circuit avoids logical erasure just as a quantum unitary circuit does.

\paragraph{Suggested replacement text.}
``Reversibility is shared by quantum circuits and by reversible classical circuits. What is specifically quantum is the use of superposition, phase, and entanglement. The Toffoli gate is important here because it is classically reversible and also appears inside universal quantum gate constructions when combined with \(H\).''

\section*{Important Clarifications, Even If Not Strict Errors}

\subsection*{6. The treatment of phase should distinguish state vectors from physical states}

The paper tracks objects like \(e^{i\theta}|1\rangle\) as distinct account states around the \(A\)-circle. That is acceptable as a representation of vectors, but not as a representation of physical pure states, because global phase is unobservable:
\[
|\psi\rangle \sim e^{i\phi}|\psi\rangle.
\]
For a single isolated qubit, \(|1\rangle\) and \(e^{i\phi}|1\rangle\) represent the same physical pure state. Relative phase becomes physically meaningful only in superposition, for example
\[
\frac{|0\rangle+|1\rangle}{\sqrt2}
\qquad \text{versus} \qquad
\frac{|0\rangle-|1\rangle}{\sqrt2}.
\]

The paper can keep the circle notation, but it should say explicitly that it is tracking vector representatives for computational bookkeeping, not distinct physical rays.

\subsection*{7. The superposition notation is useful but incomplete unless amplitude bookkeeping is formalized}

The paper says that the accounting notation records the basis-state structure, while amplitudes are ``recorded separately.'' That is fine as a hybrid notation, but it means the T-account system alone is not yet a full representation of a general qubit state.

A general qubit is
\[
|\psi\rangle=\alpha|0\rangle+\beta|1\rangle,
\qquad
|\alpha|^2+|\beta|^2=1,
\]
and the nontrivial part of quantum mechanics lives in the amplitudes \(\alpha\) and \(\beta\), not just in the basis labels. The paper should therefore describe the proposal as a combined representation:
\[
\text{accounting basis labels} + \text{externally recorded complex amplitudes}.
\]

\subsection*{8. The explanation of the ``core'' of the Toffoli circuit should be more precise}

The paper says the sequence
\[
X_{23}T_3^\dagger X_{13}T_3X_{23}T_3^\dagger X_{13}T_3
\]
``flips'' the target qubit when both controls are \(A\). That is not literally the cleanest description. The surrounding Hadamards convert a phase construction into a bit flip:
\[
H X H = Z,
\qquad
H Z H = X.
\]
More precisely, the middle \(T\)- and controlled-\(X\)-sequence is a controlled phase construction, and conjugation by \(H_3\) turns the phase action on qubit 3 into the desired \(X\)-action on qubit 3.

\paragraph{Suggested replacement text.}
``The middle \(T\)- and controlled-\(X\)-sequence implements the phase gadget at the heart of the Toffoli construction. Conjugating the target qubit by \(H_3\) converts that phase action into the desired controlled-controlled bit flip.''

\section*{Results I Checked and Agree With}

\begin{itemize}
\item The computational basis vectors are correct:
\[
|0\rangle=\begin{bmatrix}1\\0\end{bmatrix},
\qquad
|1\rangle=\begin{bmatrix}0\\1\end{bmatrix}.
\]

\item The basic gate actions are correct:
\[
X|0\rangle=|1\rangle, \qquad X|1\rangle=|0\rangle,
\]
\[
Z|0\rangle=|0\rangle, \qquad Z|1\rangle=-|1\rangle,
\]
\[
T|1\rangle=e^{i\pi/4}|1\rangle, \qquad S|1\rangle=e^{i\pi/2}|1\rangle=i|1\rangle.
\]

\item The identities
\[
T^2=S, \qquad S^2=Z, \qquad T^4=Z
\]
are correct.

\item The phase arithmetic used in the \(A\)-circle examples is correct, for example
\[
S\!\left(e^{i3\pi/4}|1\rangle\right)=e^{i5\pi/4}|1\rangle=-e^{i\pi/4}|1\rangle.
\]

\item The Hadamard action is correct:
\[
H|0\rangle=\frac{|0\rangle+|1\rangle}{\sqrt2},
\qquad
H|1\rangle=\frac{|0\rangle-|1\rangle}{\sqrt2}.
\]

\item The recomposition idea is mathematically legitimate because scalar phase can be shifted across tensor factors:
\[
e^{i\phi}|a\rangle\otimes|b\rangle
=
|a\rangle\otimes e^{i\phi}|b\rangle.
\]

\item The measurement formulas are correct:
\[
\Pr(0)=|\alpha|^2, \qquad \Pr(1)=|\beta|^2
\]
for \(\alpha|0\rangle+\beta|1\rangle\), and more generally
\[
p(m)=\langle\psi|M_m^\dagger M_m|\psi\rangle.
\]

\item The Toffoli decomposition is correct. I checked its action on all basis states:
\[
|000\rangle\mapsto|000\rangle,\;
|001\rangle\mapsto|001\rangle,\;
|010\rangle\mapsto|010\rangle,\;
|011\rangle\mapsto|011\rangle,\;
|100\rangle\mapsto|100\rangle,\;
|101\rangle\mapsto|101\rangle,
\]
\[
|110\rangle\mapsto|111\rangle,
\qquad
|111\rangle\mapsto|110\rangle.
\]
\end{itemize}

\section*{Overall Assessment}

The paper should be revised, not discarded. The concrete gate algebra is mostly sound. The main task is to scale back a few claims and correct the foundational statements so that the paper says exactly what the mathematics actually proves.

\end{document}
